% ---
% titre: Températures moyennes en Suisse (diagramme cartésien)
% theme: FA
% chapitre: Diagrammes
% sous_chapitre: Lire, interpréter et réaliser
% niveau: [10e]
% type: EX
% tags: [c-diagramme-cartesien,c-construction-diagramme,c-proportionnalite,f-individuel,o-regle,p-entrainement]
% difficulte: 1
% source: latex
% ---

\header{Températures moyennes mensuelles}


 Voici les températures moyennes mensuelles en Suisse en 2025.

\vspace{10pt}\def\arraystretch{2}
\begin{tabularx}{\linewidth}{|r*{12}{|>{\centering\arraybackslash}X}|}
\hline
\rowcolor{gray!20}\textbf{Mois} & \textbf{J} & \textbf{F} & \textbf{M} & \textbf{A} & \textbf{M} & \textbf{J} & \textbf{J} & \textbf{A} & \textbf{S} & \textbf{O} & \textbf{N} & \textbf{D}\\
\hline
\textbf{Temp. (\textdegree C)} & $-0,8$ & $0,0$ & $3,9$ & $6,6$ & $8,1$ & $16,1$ & $15,2$ & $16,4$ & $11,3$ & $7,4$ & $2,1$ & $-1,4$\\
\hline
\end{tabularx}
\def\arraystretch{1}

\vspace{20pt}
\textbf{Ce diagramme cartésien représente la température moyenne de chaque mois. Place précisément les points manquants, puis relie-les.}

\vfill
\begin{center}
\begin{tikzpicture}[x=1.05cm, y=0.45cm]
  % lignes principales (tous les 5 °C uniquement)
  \foreach \y in {-5,0,5,10,15,20}{
    \draw[gray!70] (0,\y) -- (12,\y);
  }
  % cadre et axes
  \draw[thick] (0,-5) rectangle (12,20);
  \draw[thick, ->] (0,0) -- (12.4,0);
  \draw[thick, ->] (0,-5) -- (0,21);
  % graduations de l'axe des ordonnées
  \foreach \y in {-5,0,5,10,15,20}{
    \node[left] at (0,\y) {$\y$};
  }
  \node[above] at (0,21) {Température (\textdegree C)};
  % mois en abscisse
  \foreach \x/\m in {1/Janvier,2/Février,3/Mars,4/Avril,5/Mai,6/Juin,7/Juillet,8/Août,9/Septembre,10/Octobre,11/Novembre,12/Décembre}{
    \node[rotate=90, anchor=east] at (\x-0.5,-5.4) {\m};
    \draw (\x-0.5,-5) -- (\x-0.5,-5.3);
  }
  % correction : croix en arrière-plan
  \begin{pgfonlayer}{background}
    \begin{scope}[color=\colSolution]
      \foreach \x/\t in {2/0.0,3/3.9,4/6.6,5/8.1,6/16.1,7/15.2,9/11.3,10/7.4,11/2.1,12/-1.4}{
        \node at (\x-0.5,\t) {\Large $\times$};
      }
    \end{scope}
  \end{pgfonlayer}
  % correction : courbe
  \begin{scope}[color=\colSolution]
    \draw[thick] (0.5,-0.8) -- (1.5,0.0) -- (2.5,3.9) -- (3.5,6.6) -- (4.5,8.1) -- (5.5,16.1) -- (6.5,15.2) -- (7.5,16.4) -- (8.5,11.3) -- (9.5,7.4) -- (10.5,2.1) -- (11.5,-1.4);
  \end{scope}
\end{tikzpicture}
\end{center}
\vfill
\newpage
\begin{tikzpicture}\draw[step=0.5cm, gray, very thin] (0,0) grid (\textwidth,20);\end{tikzpicture}